The flexibility of conversion dates has a profound impact on the value of a convertible bond. Table~\ref{tab:bond_valuation} summarizes the bond prices obtained under different assumptions regarding the timing of conversion opportunities.\\

Allowing conversion at any time (American-style) produces the highest bond value. This increased value reflects the maximum advantage that flexibility can provide, as bondholders have the freedom to convert at the most favorable moment. Specifically, American-style conversion achieves a premium of approximately 5.12\% over the standard scenario where conversion is permitted only at coupon dates.\\

Introducing monthly conversion opportunities captures a large proportion of this flexibility premium. Although not fully equivalent to continuous conversion, monthly opportunities allow bondholders to respond more effectively to favorable movements in the underlying stock price. This suggests that frequent discrete conversion dates provide a good approximation of continuous flexibility without the complexity of fully continuous decision-making.\\

Conversely, allowing conversion only at specific arbitrary dates can reduce the bond’s value. When conversion dates are misaligned with periods when the option to convert would be most valuable, the bondholder's ability to maximize the benefit of conversion is compromised. This demonstrates that not only the number of conversion opportunities but also their timing plays a crucial role in determining the overall bond value.\\

Several factors influence the extent to which conversion flexibility adds value. In high-volatility environments, the benefit of frequent conversion opportunities becomes more pronounced. Greater stock price variability increases the likelihood of the stock price reaching favorable levels for conversion, thereby enhancing the option’s value. When the stock price is persistently low relative to the conversion threshold, the advantage of additional flexibility diminishes, as conversion remains unlikely. Moreover, the value of flexibility tends to be greater earlier in the bond's life, when there is more time for favorable conditions to arise. As the bond approaches maturity, the time value of the conversion option declines, reducing the incremental benefit of additional opportunities.\\

The implications for both issuers and investors are significant. Issuers may choose to restrict conversion opportunities to specific dates, such as coupon payment dates, to lower the value of the conversion option and, in turn, reduce the bond’s cost. \\

For investors, greater conversion flexibility enhances the attractiveness of the bond, particularly in volatile market conditions. Bonds offering more frequent conversion rights tend to command higher prices. Evaluating the structure of conversion rights is therefore critical for investors aiming to maximize returns from convertible bond investments.\\